% Statement, parameters, logical structure (v3 §0)
\section{Statement, parameters and logical structure}\label{sec:setup}

Throughout, $L:=\log x$ and $\ell:=\log\log x$ at the scale $x$ under discussion, and $\Lambda:=\log X_0$. The exponent
$\theta\in[\tfrac14,\theta_1]$, $\theta_1:=\tfrac12-10^{-4}$, is fixed. Numerical values are quoted for $\theta=\tfrac14$, which is the
worst case for the absorption estimates of Section~\ref{sec:queue}.

\subsection{The main theorem}
\begin{theorem}[Theorem A]\label{thm:A}
% v3: Theorem A (§0.1)
Fix $\theta\in[\tfrac14,\theta_1]$, $\theta_1:=\tfrac12-10^{-4}$, and the parameters of Section~\ref{sec:params}; in particular
$X_0\ge\max(e^{10^8},X_{\mathrm{Hux}})$, where $X_{\mathrm{Hux}}$ is defined in \eqref{eq:huxley}. Let $\PP$ be the random discrete
Beurling system of Section~\ref{sec:construction}: the hard greedy rule with box-centred jitter, for the target $T$ of
\eqref{eq:target}. Then on the event $\mathcal G=\{\tau=\infty\}$ of Definition~\ref{def:good}, which has probability at least
$1-1.5/\log X_0>0$:
\begin{enumerate}[label=(\roman*)]
\item $|N_{\PP}(x)-x|\le C_{X_0}x^{\theta}+3x^{\theta}\log^*x+x^{\eta}$ for all $x\ge1$, where $\log^*$ is the iterated logarithm, $\eta:=\min(\theta,\tfrac12-\theta)/2$ and
$C_{X_0}<\infty$;
\item $\zeta_{\PP}$ has zeros $\rho_k$ $(k\ge1)$ with
$\operatorname{Re}\rho_k=1-\bigl(3\log\log x_k+O_{X_0}(1)\bigr)/\log x_k\to1$.
\end{enumerate}
\end{theorem}

The threshold $X_{\mathrm{Hux}}$ is finite by Heath-Brown's theorem \cite{HeathBrown1988} and depends on no parameter of the
construction. Its value is never used: every quantitative estimate in the proof holds for all $\log X_0\ge10^8$, with the
estimates that rest on \eqref{eq:huxley} (Lemma~\ref{lem:margin}(a) and its consumers) needing only $X_0\ge X_{\mathrm{Hux}}$ in
addition.

\begin{corollary}\label{cor:notQ}
% v3: Consequence (§0.1)
The answer to Question~\ref{def:Q} is negative. Equivalently, a discrete $[1,\beta]$-system with $\beta<\tfrac12$ exists.
\end{corollary}

\begin{proof}
Fix any $\omega\in\mathcal G$. The realised prime sequence $\PP(\omega)$ is a fixed discrete Beurling system: it has at most one
prime per slot, hence finitely many below any bound. By Theorem~\ref{thm:A}(i) it has $N=x+O(x^{\theta}\log^*x)$, and by Theorem~\ref{thm:A}(ii) it has no zero-free half-plane. That $T$ was constructed adaptively
plays no role in this conclusion. By Remark~\ref{rem:equiv}, direction (b)$\Rightarrow$(a), $\psi_{\PP}(x)-x$ has
no power saving, so $\PP(\omega)$ is a $[1,\beta]$-system with $\beta\le\theta<\tfrac12$.
\end{proof}

\begin{proof}[Proof of Theorem~\ref{thm:A}, assuming the results of Sections~\ref{sec:construction}--\ref{sec:baserange}]
Part (i) follows from the target error \eqref{eq:T1}, the property \eqref{eq:La} and $|R|\le2x^{\theta}\log^*x$. Property
\eqref{eq:La} follows from the two-sided lag \eqref{eq:lag} and the queue bound, Theorem~\ref{thm:queue}, on $\mathcal G$. The queue
bound needs two inputs.
\begin{itemize}
\item The drift of Section~\ref{sec:noise}: the noise response, the margin, and R1$'$ (Corollary~\ref{cor:r1}). The jump of $T'$
at $X_0$ is Lemma~\ref{lem:xb}, with Lemma~\ref{lem:si} in its largest range.
\item The probabilistic inputs of Section~\ref{sec:queue}: Freedman's inequality, the cluster event (Proposition~\ref{prop:ecl}) via
the rank-full counting of Appendix~\ref{app:rankfull}, the local count (Corollary~\ref{cor:localcount}), and the sparsity of the
absorbed set (Lemma~\ref{lem:lemmaP}).
\end{itemize}
Part (ii) is Theorem~\ref{thm:zeros}, from \eqref{eq:La} and Lemma~\ref{lem:ZA}. Finally, $\Prob(\mathcal G)\ge1-1.5/\log X_0$ is
Lemma~\ref{lem:budget}.
\end{proof}

\subsection{Parameters}\label{sec:params}
All parameters are deterministic functions of the scale. They are fixed here, before the construction starts.
\begin{center}\small
\begin{tabular}{@{}p{2.1cm}p{12.4cm}@{}}
\hline
parameter & value \\ \hline
$\theta$, $\theta'$ & $\theta\in[\tfrac14,\theta_1]$, $\theta_1:=\tfrac12-10^{-4}$, fixed ($\theta=\tfrac14$ suffices for Corollary~\ref{cor:notQ});
the local prime number theorem uses the exponent $\theta':=\theta+0.01$. \\
$X_0$ & $X_0\ge\max(e^{10^8},X_{\mathrm{Hux}})$: the single regime for the whole proof. $X_{\mathrm{Hux}}$ is defined in
\eqref{eq:huxley}, depends on no parameter, and gives $\mathrm{MA}[\mathbb P](q)\log q\ge0.99$ for every $q\ge X_0$
(Lemma~\ref{lem:margin}(a)). The local prime number theorem is used only above $X_{\mathrm{LP}}(X_0)$, with
$\log X_{\mathrm{LP}}=O(\Lambda^2\log^2\Lambda)$ \eqref{eq:XLP}. $X_{\mathrm{Hux}}$ is effective but has not been computed, so $X_0$
is not fixed at $e^{10^8}$; every quantitative estimate therefore holds uniformly in $\Lambda=\log X_0\ge10^8$, in particular the
computer-assisted Theorem~\ref{thm:bs}, Theorem~\ref{thm:caplw} and \eqref{eq:mu1sum}. \\
$\kappa$ & $\kappa:=X_0^{-5/12+2\eps}$ with $\eps:=0.01$, i.e.\ $\kappa=X_0^{-0.3967}\approx e^{-3.97\cdot10^7}$ (the window of the
rigid density; every use needs only $0<\eps\le0.04$). \\
$H(x)$ & $H:=\max(H_{\mathrm{std}},M_1^2)$, $H_{\mathrm{std}}(x):=\exp(8\ell^2)$, so $H(X_0)\approx e^{2715}$. The raise $M_1^2$ can
exceed $H_{\mathrm{std}}$ only on the raised range of Theorem~\ref{thm:caplw}, which ends at
$\ell_{\mathrm{std}}=\sqrt{(7.5\Lambda+\log4)/4}$ ($1.37\cdot10^4$ at $\Lambda=10^8$). Always $\log H-\log H_{\mathrm{std}}\le10^{-6}L$,
so $H=x^{o(1)}$; also $H\le16X_0^{15}$ for $\ell\le\ell_{\mathrm{std}}$, and $H=H_{\mathrm{std}}$ beyond. \\
$\lambda$ & $\lambda=4L$ (Theorem~\ref{thm:queue}); $\lambda_u=3\log u+2\log\log u$ (per-scale noise);
$\lambda'_u=\lambda_u+2+\log H(u)$ (short windows); $\lambda_P=4\log P$ (Lemma~\ref{lem:landing}). \\
$c'$ & $c'=0.15$ (the drift margin). \\
$r_0$, $c$, $M_1$ & $r_0=\lceil L/(3\log L)\rceil+3$. Capacity threshold $\log c=\max(4\ell^2-\log4,\gamma(s)\Lambda)$,
$s:=\ell-2\log\Lambda$, with the explicit ramp $\gamma\le7.5$ of Theorem~\ref{thm:caplw}; so $c=\tfrac14H_{\mathrm{std}}^{1/2}$ for
$\ell\ge\ell_{\mathrm{std}}$, and $c\le X_0^{7.5}$ for $\ell\le\ell_{\mathrm{std}}$. $M_1=r_0+\varphi(r_0)+1$ with $\varphi(u)=c+uc^{1/2}+u(u+1)/2$, so
$M_1\le2c$. Tolerance: $M_1\le H/L^{14}$. \\
$L_0$, $q_{\min}$, $B$ & $L_0=\max(L^5,48M_1L^3/c'^2)$; $q_{\min}=L_0L^6$; queue bound $B(x)\le5M_1L^2/c'$ \eqref{eq:queueB};
$K_{\mathrm{loc}}=2B(4P)+7\lambda_P/3+23$. \\
absorption & $L_*=10^7$ (inactive, since $\log X_0\ge10^8$); $y(n)=H(n)L_0(n)(\log n)^6$, with $y\le X_0^{22.6}L^9$ for $\ell\le\ell_{\mathrm{std}}$ and $\log2y\le12.1\ell^2$ beyond; $\Abs=\emptyset$ wherever
$y<\Xm=X_0-H(X_0)/2$, which holds for $\ell\le2\log\Lambda+s_H$, $s_H=5.9745$ (Theorem~\ref{thm:caplw}(b)); $h(u)=u^{\theta/2}$;
Lemma~\ref{lem:lemmaP} uses $\sup_{[u,2u]}y$. \\
cluster window & $\delta_1=2L^{-6}$. \\
\eqref{eq:La} & $|D|\le x^{\eta}$ with $\eta=\min(\theta,\tfrac12-\theta)/2$, for all $x\ge\Xm$; $D=0$ below $\Xm$. \\
amplitudes & $\kappa'_k=\kappa_0/\log^2x_k$, $\kappa_0\in(0,\tfrac14]$; $\Gamma_k=\kappa_0x_k^{1-\theta}/(2\log^3x_k)$;
$\gamma_k\in[\Gamma_k,2\Gamma_k]$; $\alpha_k\in[1,2]$; $\kappa_+=0.8$. \\
switch-on & $x_1\ge\exp(X_0^2)$, $x_{k+1}\ge\exp(x_k)$; $R_k\equiv0$ on $[1,x_k+2H(x_k)]$. \\
budget & $\Prob(\tau<\infty)\le\sum_m(\Lambda+m\log2)^{-2}+O(X_0^{-1})\le1.5/\log X_0$ (Lemma~\ref{lem:budget}). \\ \hline
\end{tabular}
\end{center}
This regime is far beyond anything that can be simulated. Experiments at small scales (with $X_0\le10^7$ and
$x\le5\cdot10^8$) reproduce the qualitative behaviour (an observed queue length of order $\log^2x$, far below the proved bound
$B=e^{O(\ell^2)}$, and designed zeros near the predicted points), but they test only the mechanism, not any inequality used in
the proof.

\subsection{Logical structure}\label{sec:logic}
The table lists the main results with their inputs. An input marked $\downarrow$ is used only at a strictly earlier time or a
strictly smaller height; every other input is used at the same time. The same-time dependencies form an acyclic graph; this was
checked mechanically on the full dependency graph (79 nodes, 194 edges, 28 of them $\downarrow$; Appendix~\ref{app:cert-dep}).
The table is an abridgement of that graph. The
time-ordering of the $\downarrow$ edges is Section~\ref{sec:noncirc}.
\begin{center}\footnotesize
\begin{tabular}{@{}p{4.4cm}p{6.6cm}p{3.5cm}@{}}
\hline
result & inputs & used in \\ \hline
\eqref{eq:huxley} & Heath-Brown, de la Vall\'ee Poussin & \eqref{eq:G1}, Lemmas~\ref{lem:blin}, \ref{lem:margin} \\
Lemma~\ref{lem:adapted} (adaptedness) & Lemma~\ref{lem:C2} (causality), Section~\ref{sec:params} & Section~\ref{sec:noncirc}, Lemma~\ref{lem:budget} \\
\eqref{eq:G1} (positivity) & \eqref{eq:huxley}, the cap, \eqref{eq:ampl} & slots, \ref{eq:sF1} \\
\eqref{eq:T1} (target error) & Rankin; Lemma~\ref{lem:lemmaP}$\downarrow$ & Theorem~\ref{thm:A}(i), Section~\ref{sec:noise}, Lemma~\ref{lem:ZA} \\
Lemmas~\ref{lem:count}--\ref{lem:excess} & the greedy rule & Theorem~\ref{thm:queue}, Corollary~\ref{cor:localcount}, \eqref{eq:La} \\
Lemmas~\ref{lem:ZB}, \ref{lem:ZBlow} & \eqref{eq:rhodet}; partial summation & Section~\ref{sec:zeros}; Theorem~\ref{thm:lp} \\
Lemma~\ref{lem:ZA} & Lemma~\ref{lem:lemmaP}$\downarrow$, Theorem~\ref{thm:caplw}(b)$\downarrow$ & Lemma~\ref{lem:goodheights}, Theorem~\ref{thm:zeros} \\
Lemmas~\ref{lem:designedzero}--\ref{lem:goodheights} & Lemma~\ref{lem:ZB}; mean values; \eqref{eq:La}$\downarrow$ on $[1,x_k)$, Lemma~\ref{lem:ZA}$\downarrow$ & Theorem~\ref{thm:zeros}, choice of $\gamma_k$ \\
\eqref{eq:La} & \eqref{eq:lag}, Theorem~\ref{thm:queue}, Lemma~\ref{lem:budget} & Theorem~\ref{thm:A}, Theorem~\ref{thm:zeros} \\
Theorem~\ref{thm:zeros} & \eqref{eq:La}, Lemmas~\ref{lem:ZB}, \ref{lem:ZA}, \ref{lem:designedzero}--\ref{lem:goodheights}; Rouch\'e & Theorem~\ref{thm:A}(ii) \\
Lemmas~\ref{lem:blin}, \ref{lem:margin} & \eqref{eq:huxley}; Lemma~\ref{lem:cm}$\downarrow$, \eqref{eq:T1}$\downarrow$, \eqref{eq:La}$\downarrow$ on cofactors & Proposition~\ref{prop:g1} \\
Lemma~\ref{lem:cm} (Chebyshev--Mertens) & \eqref{eq:HBS} from Theorem~\ref{thm:queue}$\downarrow$ & Section~\ref{sec:noise}, \ref{eq:sF3p}, \ref{eq:sF4} \\
Theorem~\ref{thm:bs}, Lemma~\ref{lem:bs-profile} & Lemmas~\ref{lem:cm}, \ref{lem:smooth-v}; certificate & Corollary~\ref{cor:saias}, Lemma~\ref{lem:si} \\
Theorems~\ref{thm:kr}, \ref{thm:kr-random}, Corollary~\ref{cor:kr-weights}, Proposition~\ref{prop:true-kernels} & Lemmas~\ref{lem:lemmaM}, \ref{lem:tau-near}, Corollary~\ref{cor:saias}, \eqref{eq:RH}; for the absorbed count: Corollary~\ref{cor:lag}, Lemma~\ref{lem:lemmaP}, Proposition~\ref{prop:ecl} & Proposition~\ref{prop:br}, Lemma~\ref{lem:xb} \\
Lemma~\ref{lem:threshold-mass} (threshold mass) & Lemma~\ref{lem:cm}; Corollary~\ref{cor:localcount}$\downarrow$ & Proposition~\ref{prop:g1} \\
Lemma~\ref{lem:noise} (per-scale noise) & Freedman; Theorem~\ref{thm:queue}$\downarrow$, Proposition~\ref{prop:ecl}$\downarrow$, Corollary~\ref{cor:localcount}$\downarrow$, Lemma~\ref{lem:lemmaP}$\downarrow$ & Proposition~\ref{prop:pointwise} \\
Theorem~\ref{thm:lp} (local PNT) & Lemma~\ref{lem:ZBlow}, \eqref{eq:Rhatbound}, \eqref{eq:T1}, Lemma~\ref{lem:cm}; $|N-T_0|$ bound$\downarrow$ & Proposition~\ref{prop:g1}, Lemma~\ref{lem:si} \\
\eqref{eq:osc} & \eqref{eq:rhodet}, \eqref{eq:target}; Lemma~\ref{lem:rhoA}(c) via Lemma~\ref{lem:lemmaP}$\downarrow$ & Propositions~\ref{prop:br}, \ref{prop:pointwise} \\
Lemmas~\ref{lem:xb}, \ref{lem:si} (jump at $X_0$) & Theorem~\ref{thm:kr-random}, Corollary~\ref{cor:kr-weights}, Propositions~\ref{prop:true-kernels}, \ref{prop:br}, Lemma~\ref{lem:tau-near}, Theorem~\ref{thm:lp}, Lemma~\ref{lem:bs-profile}; Corollary~\ref{cor:localcount}$\downarrow$, Lemma~\ref{lem:lemmaP}$\downarrow$ & Propositions~\ref{prop:g1}, \ref{prop:pointwise}, Theorem~\ref{thm:g2} \\
Proposition~\ref{prop:g1}, Theorem~\ref{thm:g2}, Proposition~\ref{prop:pointwise} & Lemma~\ref{lem:margin}, Proposition~\ref{prop:br}, Theorem~\ref{thm:lp}, Lemmas~\ref{lem:xb}, \ref{lem:noise}, \eqref{eq:osc} & Corollary~\ref{cor:r1} \\
Corollary~\ref{cor:r1} (R1$'$, every window) & Proposition~\ref{prop:pointwise} & Theorem~\ref{thm:queue} \\
Lemma~\ref{lem:landing}, Corollary~\ref{cor:localcount} & \ref{eq:sF1}, \ref{eq:sF2}, Freedman, Lemma~\ref{lem:count}, Theorem~\ref{thm:caplw}(c); Theorem~\ref{thm:queue}$\downarrow$ & Section~\ref{sec:noise}, Lemma~\ref{lem:lemmaP} ($\downarrow$ uses) \\
Theorem~\ref{thm:queue} (queue bound) & Corollary~\ref{cor:r1}, \ref{eq:sF1}, \ref{eq:sF2}, Lemmas~\ref{lem:lindley}, \ref{lem:excess}, Theorem~\ref{thm:caplw}(b),(c), Freedman; Proposition~\ref{prop:ecl}$\downarrow$, Corollary~\ref{cor:increment}$\downarrow$, \ref{eq:sF3p}$\downarrow$, Theorem~\ref{thm:queue}$\downarrow$ & \eqref{eq:La}, Lemma~\ref{lem:budget} \\
Definition~\ref{def:good}, Lemma~\ref{lem:budget} & Theorem~\ref{thm:queue}, Proposition~\ref{prop:ecl}, Lemmas~\ref{lem:landing}, \ref{lem:noise} & Theorem~\ref{thm:A} \\
Proposition~\ref{prop:ecl}, Corollary~\ref{cor:increment} & Theorem~\ref{thm:R}, Lemma~\ref{lem:heavy}, Theorem~\ref{thm:caplw} & Theorem~\ref{thm:queue}, Lemmas~\ref{lem:noise}, \ref{lem:lemmaP}, Proposition~\ref{prop:true-kernels} \\
Theorem~\ref{thm:caplw}(a) (capacity) & combinatorial covering lemmas, Rankin, Rosser--Schoenfeld & Proposition~\ref{prop:ecl} \\
Theorem~\ref{thm:caplw}(b),(c) & none (definitions of the parameters) & \ref{eq:sF3p}, Lemmas~\ref{lem:landing}, \ref{lem:lemmaP}, \ref{lem:ZA}, Theorem~\ref{thm:queue} \\
Lemmas~\ref{lem:lemmaP}, \ref{lem:rhoA} (absorption) & Lemma~\ref{lem:cm}, Theorem~\ref{thm:caplw}(b), Corollary~\ref{cor:increment}; Proposition~\ref{prop:ecl}$\downarrow$, Corollary~\ref{cor:localcount}$\downarrow$ & \eqref{eq:target}, \eqref{eq:T1}, Lemma~\ref{lem:ZA}, \eqref{eq:osc}, Proposition~\ref{prop:true-kernels} \\
Theorem~\ref{thm:R} (Appendix~\ref{app:rankfull}) & \ref{eq:sF1}--\ref{eq:sF4} & Proposition~\ref{prop:ecl} \\ \hline
\end{tabular}
\end{center}
The three loops of the proof close only through $\downarrow$ edges:
\begin{itemize}
\item queue $\to$ R1$'$ $\to$ per-scale noise $\to$ queue: Lemma~\ref{lem:noise} reads the queue at earlier slots;
\item queue $\to$ cluster event $\to$ Lemma~\ref{lem:cm} $\to$ \eqref{eq:HBS} $\to$ queue: the cluster event reads Mertens bounds at
heights $\le s-H$;
\item local count $\to$ queue $\to$ local count: Corollary~\ref{cor:localcount} at height $z$ is read only at heights $\ge4z$ or
$\ge z^2$.
\end{itemize}

\subsection{Why the induction is not circular}\label{sec:noncirc}
The induction runs over \emph{time} (slots), not over scales. Every bad event is frozen at a time before any consumer reads it.

\emph{(N0) Deterministic parameters.} $H$, $M_1$, $c$, $L_0$, $y$, $r_0$ and the switch-on points are deterministic functions of the
scale, fixed in Section~\ref{sec:params} before the construction starts. The only random objects are the decisions $P_j$ and the
jitters $U_j$.

\emph{(N1) First failure} (Definition~\ref{def:good}, Lemma~\ref{lem:budget}).
\begin{itemize}
\item Every bad event $E$ carries a time $t(E)$, a slot, and is $\mathcal F_{t(E)}$-measurable. Put
\[
  \Good(<t):=\{Q_j\le B(s_j)\text{ for all }j\text{ with }t_j<t\}\cap\bigcap_{t(E)<t}E,
\]
where the queue constraint at $j$ carries the time $t_j:=\lceil s_{j+1}\rceil+1$, the reach of its outer deterministic window.
Slots can be longer than $1$, but $s_{j+1}-s_j\le11.3L$ by \eqref{eq:G1}. The time $t_j$ is predictable, since $s_{j+1}$ is
determined by $T$ on $[1,s_{j+1}]$ (Lemma~\ref{lem:adapted}). Then $\tau:=\min\{t:\Good(\le t)\text{ fails}\}$ is an
$(\mathcal F_t)$-stopping time.
\item Each event's bound holds on $\Good(<t(E))$. Freedman's inequality is used in the form $\Prob(S\ge a,\ V\le v)$, or peeled,
so no conditioning on the future enters.
\item (D) The queue bound at $s_j$ is deterministic on $\Good(<t_j)\cap\bigcap_{t(E)\le t_j}E$ (Theorem~\ref{thm:queue}(ii)). The
relevant window event has time $\max I\le t_j$. Events with time in $(s_j,t_j]$ read only lower heights, so none of them uses $Q_j$.
\item Hence $\Prob(\tau<\infty)\le\sum_E\Prob(E^c\cap\Good(<t(E)))\le1.5/\log X_0$, and on $\mathcal G=\{\tau=\infty\}$ every event
holds at every scale.
\end{itemize}

\emph{(N2) The inputs of the queue bound at $s_j\in[x,2x)$ all have time $\le t_j$.} R1$'$ (Corollary~\ref{cor:r1}) on a window $I$
with $\max I\le t_j$ holds on $\Good(<\max I)$, and it reads:
\begin{itemize}
\item the local prime number theorem (Theorem~\ref{thm:lp}, against the target without the absorbed part) on $[1,u(1-2\delta)]$,
which is pathwise in the realised count below, so no event is needed;
\item local counts at heights $\le u/4$ or $\le\sqrt u$, and one local count at height $u(1-\delta)$
(the proof of Proposition~\ref{prop:g1}(B), Step~1), read before $u-\delta u+7H<u$;
\item the per-scale noise event for the noise on $[1,u/2]$, at time $2v\le u$;
\item the short-window events, internally ordered by window start, each window's variance bound coming from strictly earlier
windows (Proposition~\ref{prop:pointwise});
\item the cluster event at heights $\le8u/L^{11}+H$.
\end{itemize}

\emph{(N3) The cluster event follows from Theorem~\ref{thm:R} at lower heights, independently of the queue.}
\begin{itemize}
\item The moment bounds of Section~\ref{sec:queue} and Appendix~\ref{app:rankfull} carry the good-event indicator progressively,
$\mathbf 1(\Good(<s_{p_1}))\ge\mathbf 1(\Good(<s_{p_2}))\ge\cdots$ at successive pivots.
\item They use it only through: the harmonic sums \ref{eq:sF3p} over elements whose primes are all decided before $s$; Mertens
bounds at heights $\le s-H$, which hold surely on $\Good(<s)$ (Lemma~\ref{lem:cm}, local form); and the sure slot facts
\ref{eq:sF1}, \ref{eq:sF2}.
\item No queue bound at heights in the window is used quantitatively, and the window's time is at most the cluster event's time.
\end{itemize}

\emph{(N4) The local count uses the queue bound strictly below the current scale.}
\begin{itemize}
\item The local count at height $z$ (Corollary~\ref{cor:localcount}: Lemma~\ref{lem:landing} plus the count identity) reads queue
bounds at slots $\le z+H(z)$ and the landing event at height $z$.
\item Its uses are all at heights $\le$ (the consumer's height)$/4$ or $\le\sqrt{\text{consumer's height}}$: the two steps of
Lemma~\ref{lem:threshold-mass}, the bridge and near-threshold terms (Proposition~\ref{prop:br}, Lemma~\ref{lem:tau-near}),
Lemma~\ref{lem:noise}, and optionally Lemma~\ref{lem:lemmaP}.
\item The per-prime increment in Theorem~\ref{thm:queue} uses the cluster event (Corollary~\ref{cor:increment}), never the local
count, because a local count at height $x/q$ would force $B(x)\ge(3/c')L^2B(x/q)$.
\end{itemize}

\emph{(N5) Absorption and zeros.}
\begin{itemize}
\item Membership in $\Abs$ is decided at the latest prime $r<y(n)\ll n$, and $\rho_{\Abs}$ at $u$ reads only data decided before
$u-H$ (Lemma~\ref{lem:adapted}).
\item Lemma~\ref{lem:lemmaP} at $u$ reads the cluster event and local counts at heights $u/d<u$: an induction on height.
\item $\gamma_k$ is chosen from $Z_\infty$ (deterministic) and the system on $[1,x_k)$ (Lemma~\ref{lem:goodheights}), and
$R_k\equiv0$ on $[1,x_k+2H(x_k)]$. So $T$ is adapted.
\item \eqref{eq:La} and Section~\ref{sec:zeros} are pathwise consequences on $\mathcal G$.
\end{itemize}
